\begin{explanationbox}
	\textbf{\textcolor{CExplanation}{先想一道题：只知道 $5,6,7$，高在哪里？}}

\THeronFigure{5}{6}{7}{1}{0}{5}{6}{7}{已知三边 $5,6,7$；高和角都没有给出，直接用海伦公式求 $S$。}
	图中三边已足够确定面积；与其另设高或先求角，不如直接使用已有边长。

	\medskip
\par\noindent\begin{minipage}{\linewidth}
	\textbf{\textcolor{CExplanation}{第一眼应该看什么？}}
	\begin{itemize}[leftmargin=*,itemsep=0.35em]
		\item 三条边\textbf{已经知道}，但高或夹角\textbf{不知道}：优先考虑海伦公式。
		\item 已知\textbf{两边和它们的夹角}：先考虑 $S=\tfrac12 ab\sin C$，不必强求第三边。
		\item 已知\textbf{直角边}或等腰三角形的高\textbf{很容易作出}：先用 $S=\tfrac12\times\text{底}\times\text{高}$。
		\item 已知\textbf{面积、周长，并含未知边}：把海伦公式平方，建立一元二次方程，可能反求未知边。
	\end{itemize}
\end{minipage}\par

	\textbf{\textcolor{CExplanation}{为什么要用半周长 $p$？}}
	\[
		p-a=\frac{b+c-a}{2},\quad p-b=\frac{c+a-b}{2},\quad p-c=\frac{a+b-c}{2}.
	\]
	这些量都是正数：例如 $p-a>0$，就来自 ``$b+c>a$''。所以根号里的四个数都为正，面积算出来才合理。

	\textbf{\textcolor{CExplanation}{真正值得理解的来源：把未知的高消掉}}
	以一条较长的边作底，作出高 $h$，垂足把底边分成两段。
	对高左右的两个直角三角形分别使用\textbf{勾股定理}；把两式相减，$h^2$ 会消失。
	再利用\textbf{平方差公式}，面积平方就能拆成四个只含三边的因式。

	\begin{center}
		% True 13-14-15 triangle, longest side AB=c.
% T037: exact side-length model, A=(0,0), B=(c,0).
% Arguments a=BC,b=CA,c=AB,scale,mode,BC label,CA label,AB label,caption.
% Modes 0 area, 1 common altitude, 2 tangent pairs, 3 angle, 4 altitude example, 5 area pieces.
\providecommand{\THeronFigure}[9]{%
\begin{center}
\begin{tikzpicture}[scale=#4,font=\small,line join=round,line cap=round]
\pgfmathsetmacro{\tq}{((#2)^2+(#3)^2-(#1)^2)/(2*(#3))}
\pgfmathsetmacro{\th}{sqrt((#2)^2-\tq^2)}
\pgfmathsetmacro{\tp}{((#1)+(#2)+(#3))/2}
\pgfmathsetmacro{\tr}{(#3)*\th/(2*\tp)}
\pgfmathsetmacro{\tix}{((#2)*(#3)+(#3)*\tq)/(2*\tp)}
\coordinate (A) at (0,0);\coordinate (B) at (#3,0);
\coordinate (C) at (\tq,\th);\coordinate (H) at (\tq,0);
\coordinate (I) at (\tix,\tr);
\coordinate (D) at ($(B)!(I)!(C)$);
\coordinate (E) at ($(C)!(I)!(A)$);
\coordinate (F) at ($(A)!(I)!(B)$);
\ifnum#5=0
 \fill[blue!10] (A)--(B)--(C)--cycle;
 \node at ({(#3+\tq)/3},{\th/3}) {$S$};
\fi
\ifnum#5=7
 \fill[blue!10] (A)--(B)--(C)--cycle;
 \pic[draw,angle radius=3mm] {right angle=A--C--B};
\fi
\ifnum#5=1
 \fill[blue!13] (A)--(H)--(C)--cycle;
 \fill[orange!18] (B)--(H)--(C)--cycle;
\fi
\ifnum#5=5
 \fill[blue!13] (I)--(B)--(C)--cycle;
 \fill[orange!18] (I)--(C)--(A)--cycle;
 \fill[green!15!white] (I)--(A)--(B)--cycle;
\fi
\draw[thick] (A)--(B)--(C)--cycle;
\ifnum#5=2\else
 \path (B)--node[pos=.82,right=5pt] {$#6$} (C);
 \path (C)--node[pos=.25,left=5pt] {$#7$} (A);
 \ifnum#5=1\else
  \ifnum#5=4
   \path (A)--node[midway,below=28pt] {$#8$} (B);
  \else
   \ifnum#5=5
    \path (A)--node[midway,below=28pt] {$#8$} (B);
   \else
    \path (A)--node[midway,below=12pt] {$#8$} (B);
   \fi
  \fi
 \fi
\fi
\foreach \v/\where in {A/below left,B/below right,C/above}
 {\fill (\v) circle[radius=1.4pt];\node[\where=4pt] at (\v) {$\v$};}
\ifnum#5=1
 \draw[very thick,CProof] (C)--node[midway,right=4pt] {$h$} (H);
 \pic[draw,angle radius=2.8mm] {right angle=C--H--B};
 \fill (H) circle[radius=1.2pt];\node[below=4pt] at (H) {$D$};
 \path (A)--node[midway,below=16pt] {$x$} (H);
 \path (H)--node[midway,below=16pt] {$c-x$} (B);
\fi
\ifnum#5=4
 \draw[very thick,CProof] (C)--node[midway,right=4pt] {$h$} (H);
 \pic[draw,angle radius=2.8mm] {right angle=C--H--B};
 \fill (H) circle[radius=1.2pt];\node[below=4pt] at (H) {$D$};
\fi
\ifnum#5=3
 \coordinate (K) at ($(B)!(A)!(C)$);
 \draw[CProof,dashed,thick] (A)--node[pos=.55,above=6pt] {$h_a$} (K);
 \pic[draw,angle radius=2.8mm] {right angle=A--K--C};
 \pic[draw,CProof,angle radius=9mm] {angle=A--C--B};
\fi
\ifnum#5=2\THeronIncircleDetails{2}\fi
\ifnum#5=5\THeronIncircleDetails{5}\fi
\ifnum#5=0
 \node[anchor=north,align=center,text width=11cm] at (#3/2,{-1.25/#4}) {#9};
\else
 \node[anchor=north,align=center,text width=11cm] at (#3/2,{-1.65/#4}) {#9};
\fi
\end{tikzpicture}
\end{center}}

\providecommand{\THeronIncircleDetails}[1]{%
\draw[CStatement,thick] (I) circle[radius=\tr];
\draw[gray,dashed] (I)--(A) (I)--(B) (I)--(C);
\draw[very thick,CStatement] (I)--(D) (I)--(E) (I)--node[midway,right=3pt] {$r$} (F);
\pic[draw,angle radius=2mm] {right angle=I--D--C};
\pic[draw,angle radius=2mm] {right angle=I--E--A};
\pic[draw,angle radius=2mm] {right angle=I--F--B};
\foreach \v/\where in {D/above right,E/above left,F/below}
 {\fill (\v) circle[radius=1.3pt];\node[\where=4pt] at (\v) {$\v$};}
\fill (I) circle[radius=1.6pt];\node at ($(I)+(1.3,0)$) {$I$};
\ifnum#1=2
 \draw[very thick,CStatement] (A)--node[midway,left=6pt] {$x$} (E)
  (A)--node[midway,below=16pt] {$x$} (F);
 \draw[very thick,CExample] (B)--node[midway,right=6pt] {$y$} (D)
  (B)--node[midway,below=16pt] {$y$} (F);
 \draw[very thick,CProof] (C)--node[midway,left=6pt] {$z$} (E)
  (C)--node[midway,right=6pt] {$z$} (D);
\fi}

\THeronFigure{13}{14}{15}{0.48}{1}{a}{b}{c}{底边 $AB=c$；两个直角三角形共用同一条高 $CD=h$。}

	\end{center}

\par\noindent\begin{minipage}{\linewidth}
	\textbf{\textcolor{CExplanation}{另一个令人容易记住的理解：内切圆}}
	同一顶点引出的两条切线段等长，图中同色两段对应同一个字母。
% D on BC, E on CA, F on AB, matching the formal proof.
% T037: exact side-length model, A=(0,0), B=(c,0).
% Arguments a=BC,b=CA,c=AB,scale,mode,BC label,CA label,AB label,caption.
% Modes 0 area, 1 common altitude, 2 tangent pairs, 3 angle, 4 altitude example, 5 area pieces.
\providecommand{\THeronFigure}[9]{%
\begin{center}
\begin{tikzpicture}[scale=#4,font=\small,line join=round,line cap=round]
\pgfmathsetmacro{\tq}{((#2)^2+(#3)^2-(#1)^2)/(2*(#3))}
\pgfmathsetmacro{\th}{sqrt((#2)^2-\tq^2)}
\pgfmathsetmacro{\tp}{((#1)+(#2)+(#3))/2}
\pgfmathsetmacro{\tr}{(#3)*\th/(2*\tp)}
\pgfmathsetmacro{\tix}{((#2)*(#3)+(#3)*\tq)/(2*\tp)}
\coordinate (A) at (0,0);\coordinate (B) at (#3,0);
\coordinate (C) at (\tq,\th);\coordinate (H) at (\tq,0);
\coordinate (I) at (\tix,\tr);
\coordinate (D) at ($(B)!(I)!(C)$);
\coordinate (E) at ($(C)!(I)!(A)$);
\coordinate (F) at ($(A)!(I)!(B)$);
\ifnum#5=0
 \fill[blue!10] (A)--(B)--(C)--cycle;
 \node at ({(#3+\tq)/3},{\th/3}) {$S$};
\fi
\ifnum#5=7
 \fill[blue!10] (A)--(B)--(C)--cycle;
 \pic[draw,angle radius=3mm] {right angle=A--C--B};
\fi
\ifnum#5=1
 \fill[blue!13] (A)--(H)--(C)--cycle;
 \fill[orange!18] (B)--(H)--(C)--cycle;
\fi
\ifnum#5=5
 \fill[blue!13] (I)--(B)--(C)--cycle;
 \fill[orange!18] (I)--(C)--(A)--cycle;
 \fill[green!15!white] (I)--(A)--(B)--cycle;
\fi
\draw[thick] (A)--(B)--(C)--cycle;
\ifnum#5=2\else
 \path (B)--node[pos=.82,right=5pt] {$#6$} (C);
 \path (C)--node[pos=.25,left=5pt] {$#7$} (A);
 \ifnum#5=1\else
  \ifnum#5=4
   \path (A)--node[midway,below=28pt] {$#8$} (B);
  \else
   \ifnum#5=5
    \path (A)--node[midway,below=28pt] {$#8$} (B);
   \else
    \path (A)--node[midway,below=12pt] {$#8$} (B);
   \fi
  \fi
 \fi
\fi
\foreach \v/\where in {A/below left,B/below right,C/above}
 {\fill (\v) circle[radius=1.4pt];\node[\where=4pt] at (\v) {$\v$};}
\ifnum#5=1
 \draw[very thick,CProof] (C)--node[midway,right=4pt] {$h$} (H);
 \pic[draw,angle radius=2.8mm] {right angle=C--H--B};
 \fill (H) circle[radius=1.2pt];\node[below=4pt] at (H) {$D$};
 \path (A)--node[midway,below=16pt] {$x$} (H);
 \path (H)--node[midway,below=16pt] {$c-x$} (B);
\fi
\ifnum#5=4
 \draw[very thick,CProof] (C)--node[midway,right=4pt] {$h$} (H);
 \pic[draw,angle radius=2.8mm] {right angle=C--H--B};
 \fill (H) circle[radius=1.2pt];\node[below=4pt] at (H) {$D$};
\fi
\ifnum#5=3
 \coordinate (K) at ($(B)!(A)!(C)$);
 \draw[CProof,dashed,thick] (A)--node[pos=.55,above=6pt] {$h_a$} (K);
 \pic[draw,angle radius=2.8mm] {right angle=A--K--C};
 \pic[draw,CProof,angle radius=9mm] {angle=A--C--B};
\fi
\ifnum#5=2\THeronIncircleDetails{2}\fi
\ifnum#5=5\THeronIncircleDetails{5}\fi
\ifnum#5=0
 \node[anchor=north,align=center,text width=11cm] at (#3/2,{-1.25/#4}) {#9};
\else
 \node[anchor=north,align=center,text width=11cm] at (#3/2,{-1.65/#4}) {#9};
\fi
\end{tikzpicture}
\end{center}}

\providecommand{\THeronIncircleDetails}[1]{%
\draw[CStatement,thick] (I) circle[radius=\tr];
\draw[gray,dashed] (I)--(A) (I)--(B) (I)--(C);
\draw[very thick,CStatement] (I)--(D) (I)--(E) (I)--node[midway,right=3pt] {$r$} (F);
\pic[draw,angle radius=2mm] {right angle=I--D--C};
\pic[draw,angle radius=2mm] {right angle=I--E--A};
\pic[draw,angle radius=2mm] {right angle=I--F--B};
\foreach \v/\where in {D/above right,E/above left,F/below}
 {\fill (\v) circle[radius=1.3pt];\node[\where=4pt] at (\v) {$\v$};}
\fill (I) circle[radius=1.6pt];\node at ($(I)+(1.3,0)$) {$I$};
\ifnum#1=2
 \draw[very thick,CStatement] (A)--node[midway,left=6pt] {$x$} (E)
  (A)--node[midway,below=16pt] {$x$} (F);
 \draw[very thick,CExample] (B)--node[midway,right=6pt] {$y$} (D)
  (B)--node[midway,below=16pt] {$y$} (F);
 \draw[very thick,CProof] (C)--node[midway,left=6pt] {$z$} (E)
  (C)--node[midway,right=6pt] {$z$} (D);
\fi}

\THeronFigure{13}{14}{15}{0.48}{2}{a}{b}{c}{同一顶点的两条切线段等长：$AF=AE=x$，$BF=BD=y$，$CD=CE=z$。}

\end{minipage}\par
	把三对切线段长分别记为 $x,y,z$，沿边相加：
	\[
		a=y+z,\qquad b=z+x,\qquad c=x+y,\qquad p=x+y+z.
	\]
	因而 $x=p-a$、$y=p-b$、$z=p-c$。
	内心到三边的距离都是 $r$，三块小三角形面积相加就是
	\[
		S=\tfrac12 ar+\tfrac12 br+\tfrac12 cr=pr.
	\]
	高中三角函数还能推出 $r^2=\dfrac{xyz}{x+y+z}$，于是
	\[
		S^2=p^2r^2=pxyz=p(p-a)(p-b)(p-c).
	\]
	\textbf{说明：}上面的等式不是要求初中生凭空相信；完整推导放在第三部分的第二种证明，包含每步理由。

\medskip
\par\noindent\begin{minipage}{\linewidth}
\textbf{\textcolor{CExplanation}{特殊情形一：等边三角形，海伦公式变回熟悉的公式}}

设等边三角形边长为 $\ell>0$，即 $a=b=c=\ell$。
\THeronFigure{4}{4}{4}{1.45}{0}{\ell}{\ell}{\ell}{三条边都相等，三个半周长差也相等：$p-a=p-b=p-c=\ell/2$。}
\end{minipage}\par
先代入半周长，再逐个计算根号中的四项：
\[
 p=\frac{3\ell}{2},\qquad
 p-a=p-b=p-c=\frac\ell2.
\]
于是海伦公式直接给出
\[
 \begin{aligned}
 S&=\sqrt{\frac{3\ell}{2}\cdot\frac\ell2\cdot\frac\ell2\cdot\frac\ell2}
   =\sqrt{\frac{3\ell^4}{16}}
   =\boxed{\frac{\sqrt3}{4}\ell^2}.
 \end{aligned}
\]
\textbf{回验：}等边三角形作高得 $h=\sqrt3\ell/2$，底乘高的一半也是
$\ell h/2=\sqrt3\ell^2/4$。例如边长 $6$，海伦公式算得 $S=9\sqrt3$。
这说明等边三角形面积公式正是海伦公式的一种特殊情形。

\medskip
\par\noindent\begin{minipage}{\linewidth}
\textbf{\textcolor{CExplanation}{特殊情形二：直角三角形，海伦公式变成两直角边乘积的一半}}

设 $\angle C=90^\circ$，两直角边为 $a=BC,b=CA>0$，斜边为 $c=AB$，则 $c^2=a^2+b^2$。
\THeronFigure{3}{4}{5}{1.35}{7}{a}{b}{c}{这里的 $c$ 是斜边；勾股关系 $c^2=a^2+b^2$ 是化简的关键。}
\end{minipage}\par
仍令 $p=(a+b+c)/2$。把海伦公式的四个因子配成两对：
\[
 \begin{aligned}
 p(p-c)&=\frac{(a+b+c)(a+b-c)}4
 =\frac{(a+b)^2-c^2}{4}=\frac{ab}{2},\\
 (p-a)(p-b)&=\frac{(b+c-a)(a+c-b)}4
 =\frac{c^2-(a-b)^2}{4}=\frac{ab}{2}.
 \end{aligned}
\]
两处最后一步都用到了 $c^2=a^2+b^2$。因此
\[
 S=\sqrt{p(p-c)(p-a)(p-b)}
   =\sqrt{\frac{ab}{2}\cdot\frac{ab}{2}}
   =\boxed{\frac{ab}{2}},
\]
最后一步利用 $a,b>0$，面积取正值。
例如 $3,4,5$ 直角三角形，$p=6$，直接代入为
\[
 S=\sqrt{6(6-3)(6-4)(6-5)}=\sqrt{36}=6
  =\frac12\times3\times4.
\]
\textbf{记忆要点：}海伦公式也适用于直角三角形；结果与底乘高的一半完全一致。
实际解题时，直角边已经给出就直接用 $ab/2$，这里的代入主要用于理解和验证公式。

	\textbf{\textcolor{CExplanation}{为什么说这个结论有用？}}
	\begin{itemize}[leftmargin=*,itemsep=0.35em]
		\item \textbf{少设未知数：}有三边，就无需再设高 $h$ 或角 $C$ 来求面积。
		\item \textbf{方便求其他量：}面积求出后，可求某边上的高 $h=2S/\text{该边长}$，或内切圆半径 $r=S/p$。
		\item \textbf{反过来求边：}题目给出面积和周长时，海伦公式能转成关于未知边的方程。
		\item \textbf{方便判断方法是否划算：}若需要先费力求出第三边才能套公式，它未必比常规面积公式快。
	\end{itemize}
\end{explanationbox}
